\section{Articulation Points}
\label{sec:graph-articulation-points}

\problemstatement{Given an undirected, connected graph, find every
\textbf{articulation point} (cut vertex) -- a node whose removal
(along with all its incident edges) would increase the number of
connected components.}

\begin{patternbox}
This closing problem of the entire graph series reuses
Section~\ref{sec:graph-bridges}'s exact $\textit{tin}$/$\textit{low}$
DFS, with the condition shifted from edges to \textbf{nodes}, and from
strict to non-strict inequality: a non-root node $u$ is an
articulation point if it has a child $v$ with
$\textit{low}[v] \ge \textit{tin}[u]$ -- plus one special case for the
DFS \textbf{root}, which has no parent for the usual argument to apply
to.
\end{patternbox}

\subsection*{Why $\ge$ replaces bridges' strict $>$}

A bridge requires $v$'s subtree to have \textbf{no} way back to $u$ or
higher ($\textit{low}[v] > \textit{tin}[u]$, strictly). An articulation
point is a weaker condition: $u$ is critical even if $v$'s subtree can
reach back \emph{exactly} to $u$ itself ($\textit{low}[v] =
\textit{tin}[u]$) -- because reaching back only as far as $u$ still
means removing $u$ disconnects $v$'s subtree from everything
\emph{above} $u$, even though $v$'s subtree remains internally
connected to $u$'s position. Only if $v$'s subtree can reach back
\emph{above} $u$ (to some ancestor, $\textit{low}[v] < \textit{tin}[u]$)
does removing $u$ fail to disconnect anything.

\subsection*{The root's special case}

The root of the DFS tree has no parent, so the ``can my child's
subtree reach above me'' question doesn't apply to it the same way.
Instead: the root is an articulation point exactly when it has
\textbf{more than one child} in the DFS tree -- each child starts an
independent branch that (by definition of being a separate DFS tree
child rather than reached via a back edge) has no other connection to
the rest of the graph except through the root.

\subsection*{C++ implementation}

\begin{lstlisting}
#include <bits/stdc++.h>
using namespace std;

int timer = 0;

void dfs(int node, int parent, vector<vector<int>>& adj, vector<bool>& visited,
         vector<int>& tin, vector<int>& low, vector<bool>& isArticulation) {
    visited[node] = true;
    tin[node] = low[node] = timer++;
    int children = 0;

    for (int neighbor : adj[node]) {
        if (neighbor == parent) continue;

        if (visited[neighbor]) {
            low[node] = min(low[node], tin[neighbor]);
        } else {
            dfs(neighbor, node, adj, visited, tin, low, isArticulation);
            low[node] = min(low[node], low[neighbor]);

            if (parent != -1 && low[neighbor] >= tin[node]) {
                isArticulation[node] = true;
            }
            children++;
        }
    }

    if (parent == -1 && children > 1) {
        isArticulation[node] = true;
    }
}

vector<int> findArticulationPoints(int n, vector<vector<int>>& adj) {
    vector<bool> visited(n, false), isArticulation(n, false);
    vector<int> tin(n), low(n);
    timer = 0;

    for (int i = 0; i < n; i++) {
        if (!visited[i]) dfs(i, -1, adj, visited, tin, low, isArticulation);
    }

    vector<int> result;
    for (int i = 0; i < n; i++) {
        if (isArticulation[i]) result.push_back(i);
    }
    return result;
}

int main() {
    int n = 5;
    vector<vector<int>> adj(n);
    auto addEdge = [&](int u, int v) { adj[u].push_back(v); adj[v].push_back(u); };
    addEdge(0, 1); addEdge(1, 2); addEdge(2, 0); addEdge(1, 3); addEdge(3, 4);

    for (int x : findArticulationPoints(n, adj)) cout << x << " ";
    cout << endl; // 1 3
    return 0;
}
\end{lstlisting}

\subsection*{Dry run}

\dryrun{Take edges $0{-}1,1{-}2,2{-}0,1{-}3,3{-}4$ (a triangle
$0,1,2$, with $1$ also connected via $3$ to a pendant $4$), DFS from
$0$.}

This is the same graph shape as Section~\ref{sec:graph-bridges}'s
example with an added pendant ($3{-}4$). Following the same DFS order:
$\textit{tin}/\textit{low}$ for $0,1,2$ resolve identically (the
triangle is unaffected). At node $3$ (child of $1$, with one child of
its own, $4$): $\textit{dfs}(4,3)$ gives
$\textit{tin}[4]=\textit{low}[4]=4$. Back at $3$:
$\textit{low}[3]=\min(3,4)=3$; check articulation: parent of $3$ is
$1$ ($\neq-1$), $\textit{low}[4]{=}4 \ge \textit{tin}[3]{=}3$? Yes --
\textbf{$3$ is an articulation point} (removing it disconnects $4$).
Back at $1$ (processing child $3$): $\textit{low}[1]=\min(0,
\textit{low}[3]{=}3)=0$; check articulation: parent of $1$ is $0$
($\neq-1$), $\textit{low}[3]{=}3 \ge \textit{tin}[1]{=}1$? Yes --
\textbf{$1$ is an articulation point} (removing it disconnects $\{3,4\}$
from the triangle). Node $0$ (the root): only one DFS-tree child
($1$; node $2$ was reached via a back edge, not as a tree child), so
the root special case does \textbf{not} apply -- $0$ is not an
articulation point. Final: $\{1,3\}$.

\begin{complexitybox}
\textbf{Time Complexity.} $O(V+E)$ -- a single DFS pass, identical
structure to bridge-finding.
\textbf{Space Complexity.} $O(V)$ for \textit{tin}, \textit{low},
\textit{visited}, \textit{isArticulation}, and the recursion stack.
\end{complexitybox}

\begin{takeawaybox}
This book series closes its graph coverage on a fitting note: bridges
and articulation points are, almost letter for letter, the same
algorithm, distinguished only by a strict-versus-non-strict inequality
and one root-specific special case. Across every chapter in this
series -- Greedy, Binary Search, Heap, Stack/Queue, BST, Trie,
Recursion, eight DP sub-topics, and six graph sub-topics -- the
recurring lesson has been the same: a small number of genuinely new
ideas, recognized and reapplied with small, deliberate
modifications, cover an enormous space of distinct-looking interview
problems.
\end{takeawaybox}
